\chapter{Riemannian versus Lorentzian Geometry}
\label{ch:vii30-riemannian-lorentzian}

Riemannian and Lorentzian geometry share the same differential machinery:
metrics, Levi--Civita connections, geodesics, curvature, and parallel transport.
What changes is the algebra of the metric, and with it the global geometry of
length, orthogonality, and causality.

\section{One formalism, two signatures}

A Riemannian metric has signature
\[
(n,0),
\]
while a Lorentzian metric in our convention has signature
\[
(n-1,1).
\]

\begin{proposition}[Shared local formulas]\label{prop:vii30-shared}
The Koszul formula, Christoffel symbols, geodesic equation, parallel-transport
equation, Riemann tensor, Ricci tensor, scalar curvature, and Weyl tensor have the
same coordinate formulas in both signatures.
\end{proposition}

\begin{remark}
The formulas agree because they require nondegeneracy.  Their geometric
interpretations differ because the sign of \(g(v,v)\) now matters.
\end{remark}

\section{Norm versus causal type}

In Riemannian geometry,
\[
g(v,v)>0
\qquad(v\ne0).
\]
In Lorentzian geometry a nonzero vector is timelike, null, or spacelike.

\begin{warning}[There is no single positive norm from \(g\)]\label{warn:vii30-norm}
The expression \(\sqrt{g(v,v)}\) is not a real positive norm on all Lorentzian
vectors.  Timelike vectors have negative norm-square and null vectors have zero
norm-square.
\end{warning}

\section{Proper time}

For a timelike curve \(\gamma\), define its proper-time length by
\[
\tau(\gamma)
=
\int
\sqrt{-g(\dot\gamma,\dot\gamma)}\,dt
\]
in the mostly-plus convention.

\begin{definition}[Proper time]\label{def:vii30-proper-time}
The quantity
\[
d\tau=\sqrt{-g(\dot\gamma,\dot\gamma)}\,dt
\]
is the infinitesimal proper time along a timelike curve.
\end{definition}

\begin{remark}
Riemannian geodesics locally minimize length.  Timelike Lorentzian geodesics,
under suitable local conditions and fixed endpoints, locally maximize proper time.
This reversal is one of the basic signature effects.
\end{remark}

\section{Null curves}

\begin{definition}[Null curve]\label{def:vii30-null-curve}
A regular curve is null if
\[
g(\dot\gamma,\dot\gamma)=0
\]
everywhere.
\end{definition}

Null curves have zero proper time but are geometrically nontrivial.  Null geodesics
therefore cannot be understood as shortest curves in the Riemannian sense.

\section{Time orientation}

At each point of a Lorentzian manifold, timelike vectors form two connected cones.

\begin{definition}[Time orientation]\label{def:vii30-time-orientation}
A time orientation is a continuous choice of one of the two timelike cones as
future-directed.
\end{definition}

\begin{proposition}[Criterion]\label{prop:vii30-time-orientable}
A Lorentzian manifold is time-orientable if and only if it admits a continuous
nowhere-zero timelike vector field.
\end{proposition}

\section{Causal curves}

\begin{definition}[Causal curve]\label{def:vii30-causal-curve}
A future-directed causal curve is a curve whose tangent is everywhere
future-directed timelike or null.
\end{definition}

Causal curves introduce a partial ordering flavor absent from Riemannian geometry:
one can ask whether \(q\) lies in the causal or chronological future of \(p\).

\begin{definition}[Chronological and causal future]\label{def:vii30-futures}
Write
\[
q\in I^+(p)
\]
if there exists a future-directed timelike curve from \(p\) to \(q\), and
\[
q\in J^+(p)
\]
if there exists a future-directed causal curve from \(p\) to \(q\).
\end{definition}

\section{Minkowski intervals}

For two points in Minkowski space, let
\[
\Delta x=(\Delta t,\Delta\mathbf x).
\]
Then
\[
\eta(\Delta x,\Delta x)
=
-(\Delta t)^2+|\Delta\mathbf x|^2.
\]

\begin{proposition}[Interval classification]\label{prop:vii30-interval}
The separation is timelike, null, or spacelike according as
\[
\eta(\Delta x,\Delta x)
\]
is negative, zero, or positive.
\end{proposition}

\section{Geodesics}

\begin{proposition}[Causal type is constant along a geodesic]\label{prop:vii30-geodesic-type}
For an affinely parametrized Lorentzian geodesic,
\[
g(\dot\gamma,\dot\gamma)
\]
is constant.
Therefore a nonconstant geodesic is everywhere timelike, everywhere null, or
everywhere spacelike.
\end{proposition}

\begin{proof}
Metric compatibility and
\[
\nabla_{\dot\gamma}\dot\gamma=0
\]
give
\[
\frac d{dt}g(\dot\gamma,\dot\gamma)=0.
\]
\end{proof}

\section{Orthogonality}

In Riemannian geometry,
\[
v^\perp
\]
is always a complementary positive-definite hyperplane for \(v\ne0\).  Lorentzian
orthogonality depends on causal type.

\begin{proposition}[Lorentzian orthogonal complements]\label{prop:vii30-orth}
If \(v\) is timelike, \(v^\perp\) is spacelike and positive definite.  If \(v\) is
null, then \(v\in v^\perp\), so the restriction to \(v^\perp\) is degenerate.
\end{proposition}

\section{Distance versus time separation}

Riemannian distance is symmetric:
\[
d(p,q)=d(q,p)\ge0.
\]
Lorentzian geometry instead uses causal relations and, on suitable spacetimes, a
time-separation function based on the supremum of proper times of future-directed
causal curves.

\begin{warning}[Not a metric-space distance]\label{warn:vii30-time-separation}
Lorentzian time separation is generally not symmetric and obeys a reverse
triangle-type behavior along causally ordered points.  It is not a metric in the
ordinary metric-space sense.
\end{warning}

\section{Curvature interpretation}

Positive, zero, and negative sectional curvature remain meaningful for
nondegenerate two-planes, but a Lorentzian tangent space contains planes of
different signatures.  Curvature inequalities therefore require more care than in
the Riemannian setting.

\section{Why return to Riemannian hyperbolic geometry?}

The next chapters study hyperbolic geometry, which is Riemannian and
positive-definite.  Yet the Lorentzian material is immediately useful: the
hyperboloid model realizes the hyperbolic plane inside three-dimensional Minkowski
space.

\section{Exercises}
\begin{exercise}\label{exr:vii30-01}
What signature distinguishes Riemannian from Lorentzian geometry?
\end{exercise}
\begin{hint}
Compare \((n,0)\) with index one.
\end{hint}
\begin{solution}
Riemannian signature is \((n,0)\); Lorentzian signature is \((n-1,1)\) in our convention.
\end{solution}

\begin{exercise}\label{exr:vii30-02}
Name three constructions with identical formal definitions in both signatures.
\end{exercise}
\begin{hint}
Use the Levi--Civita machinery.
\end{hint}
\begin{solution}
For example: the Levi--Civita connection, geodesic equation, and Riemann curvature tensor.
\end{solution}

\begin{exercise}\label{exr:vii30-03}
Why is \(\sqrt{g(v,v)}\) not a Lorentzian norm?
\end{exercise}
\begin{hint}
Consider timelike and null vectors.
\end{hint}
\begin{solution}
It is imaginary for timelike vectors and zero for nonzero null vectors.
\end{solution}

\begin{exercise}\label{exr:vii30-04}
Define proper time along a timelike curve.
\end{exercise}
\begin{hint}
Insert a minus sign under the square root.
\end{hint}
\begin{solution}
\(\tau=\int\sqrt{-g(\dot\gamma,\dot\gamma)}\,dt\).
\end{solution}

\begin{exercise}\label{exr:vii30-05}
What is the proper time of a null curve?
\end{exercise}
\begin{hint}
Its tangent has zero norm-square.
\end{hint}
\begin{solution}
It is zero.
\end{solution}

\begin{exercise}\label{exr:vii30-06}
Define a null curve.
\end{exercise}
\begin{hint}
Classify its tangent.
\end{hint}
\begin{solution}
It satisfies \(g(\dot\gamma,\dot\gamma)=0\) everywhere.
\end{solution}

\begin{exercise}\label{exr:vii30-07}
What is a time orientation?
\end{exercise}
\begin{hint}
Choose one timelike cone continuously.
\end{hint}
\begin{solution}
It is a continuous choice of one timelike cone as future-directed.
\end{solution}

\begin{exercise}\label{exr:vii30-08}
Give a criterion for time orientability.
\end{exercise}
\begin{hint}
Use a timelike vector field.
\end{hint}
\begin{solution}
A Lorentzian manifold is time-orientable iff it admits a continuous nowhere-zero timelike vector field.
\end{solution}

\begin{exercise}\label{exr:vii30-09}
Define a future-directed causal curve.
\end{exercise}
\begin{hint}
Allow timelike or null tangents.
\end{hint}
\begin{solution}
Its tangent is everywhere future-directed and nonspacelike.
\end{solution}

\begin{exercise}\label{exr:vii30-10}
Define \(I^+(p)\).
\end{exercise}
\begin{hint}
Use timelike curves.
\end{hint}
\begin{solution}
It is the set of points reachable from \(p\) by future-directed timelike curves.
\end{solution}

\begin{exercise}\label{exr:vii30-11}
Define \(J^+(p)\).
\end{exercise}
\begin{hint}
Allow causal curves.
\end{hint}
\begin{solution}
It is the set of points reachable from \(p\) by future-directed causal curves.
\end{solution}

\begin{exercise}\label{exr:vii30-12}
Classify a Minkowski separation with negative interval-square.
\end{exercise}
\begin{hint}
Use mostly-plus convention.
\end{hint}
\begin{solution}
It is timelike.
\end{solution}

\begin{exercise}\label{exr:vii30-13}
Classify zero interval-square.
\end{exercise}
\begin{hint}
Use the light cone.
\end{hint}
\begin{solution}
It is null.
\end{solution}

\begin{exercise}\label{exr:vii30-14}
Classify positive interval-square.
\end{exercise}
\begin{hint}
Use causal type.
\end{hint}
\begin{solution}
It is spacelike.
\end{solution}

\begin{exercise}\label{exr:vii30-15}
Why does causal type remain constant along an affine geodesic?
\end{exercise}
\begin{hint}
Differentiate \(g(\dot\gamma,\dot\gamma)\).
\end{hint}
\begin{solution}
Metric compatibility and zero covariant acceleration make it constant.
\end{solution}

\begin{exercise}\label{exr:vii30-16}
What is special about \(v^\perp\) for timelike \(v\)?
\end{exercise}
\begin{hint}
Index one leaves no negative direction.
\end{hint}
\begin{solution}
The restriction of the metric to \(v^\perp\) is positive definite.
\end{solution}

\begin{exercise}\label{exr:vii30-17}
What is special about \(v^\perp\) for null \(v\)?
\end{exercise}
\begin{hint}
Notice \(v\perp v\).
\end{hint}
\begin{solution}
It contains \(v\) in its radical, so the restricted metric is degenerate.
\end{solution}

\begin{exercise}\label{exr:vii30-18}
Is Lorentzian time separation symmetric?
\end{exercise}
\begin{hint}
Future and past differ.
\end{hint}
\begin{solution}
No, in general it is directed and not symmetric.
\end{solution}

\begin{exercise}\label{exr:vii30-19}
Does Lorentzian time separation satisfy the ordinary triangle inequality?
\end{exercise}
\begin{hint}
Think proper-time maximization.
\end{hint}
\begin{solution}
No; for causally ordered points it has a reverse triangle flavor.
\end{solution}

\begin{exercise}\label{exr:vii30-20}
What extremal behavior distinguishes short Riemannian and timelike Lorentzian geodesics?
\end{exercise}
\begin{hint}
Compare length and proper time.
\end{hint}
\begin{solution}
Riemannian geodesics locally minimize length, while timelike Lorentzian geodesics locally maximize proper time under suitable conditions.
\end{solution}

\begin{exercise}\label{exr:vii30-21}
Do null geodesics minimize ordinary distance?
\end{exercise}
\begin{hint}
There is no positive-definite distance from the Lorentzian metric.
\end{hint}
\begin{solution}
No; null geodesics are characterized by the connection and causal structure, not by Riemannian minimization.
\end{solution}

\begin{exercise}\label{exr:vii30-22}
Can sectional curvature be discussed in Lorentzian geometry?
\end{exercise}
\begin{hint}
Restrict to nondegenerate two-planes.
\end{hint}
\begin{solution}
Yes, but the signature of the plane matters and curvature inequalities require care.
\end{solution}

\begin{exercise}\label{exr:vii30-23}
Which ambient metric is used for the hyperboloid model?
\end{exercise}
\begin{hint}
Recall Minkowski space.
\end{hint}
\begin{solution}
A Lorentzian metric of signature \((2,1)\) on \(\mathbb R^3\).
\end{solution}

\begin{exercise}\label{exr:vii30-24}
State the bridge to VII/31.
\end{exercise}
\begin{hint}
Move from Lorentzian ambient space to a Riemannian submanifold.
\end{hint}
\begin{solution}
VII/31 introduces hyperbolic plane models, including the hyperboloid in Minkowski space and its disk and half-plane images.
\end{solution}

\section{Solved problem dossiers}
\begin{problem}[Shared-formulas problem]\label{prob:vii30-01}
Explain why the Levi--Civita formulas are signature independent.
\end{problem}
\begin{solution}
The Koszul and Christoffel formulas use symmetry and nondegeneracy of the metric; they do not require positivity.
\end{solution}

\begin{problem}[Proper-time problem]\label{prob:vii30-02}
Compute proper time along an inertial Minkowski line \(\gamma(t)=(t,vt)\) with \(|v|<1\).
\end{problem}
\begin{solution}
Here \(g(\dot\gamma,\dot\gamma)=-(1-|v|^2)\), so \(\tau=\sqrt{1-|v|^2}\,\Delta t\).
\end{solution}

\begin{problem}[Null-line problem]\label{prob:vii30-03}
Show \(\gamma(t)=(t,t)\) in \(1+1\) Minkowski space is null.
\end{problem}
\begin{solution}
Its tangent is \((1,1)\), whose norm-square is \(-1+1=0\).
\end{solution}

\begin{problem}[Time-orientation problem]\label{prob:vii30-04}
Show Minkowski space is time orientable.
\end{problem}
\begin{solution}
The constant vector field \(\partial_t\) is everywhere timelike. Choose the cone containing \(\partial_t\) as future-directed.
\end{solution}

\begin{problem}[Causal-future problem]\label{prob:vii30-05}
Describe \(I^+(0)\) and \(J^+(0)\) in \(1+1\) Minkowski space.
\end{problem}
\begin{solution}
With future \(t>0\), \(I^+(0)=\{t>|x|\}\) and \(J^+(0)=\{t\ge|x|,\ t\ge0\}\).
\end{solution}

\begin{problem}[Geodesic-type problem]\label{prob:vii30-06}
Prove causal type cannot change along a nonconstant affine geodesic.
\end{problem}
\begin{solution}
The squared tangent norm is constant; its sign therefore cannot change.
\end{solution}

\begin{problem}[Orthogonality problem]\label{prob:vii30-07}
Contrast the orthogonal complement of timelike and null vectors.
\end{problem}
\begin{solution}
Timelike complement is nondegenerate positive definite. A null vector belongs to its own orthogonal complement, making that restriction degenerate.
\end{solution}

\begin{problem}[Riemannian-distance problem]\label{prob:vii30-08}
Why does infimizing Lorentzian proper time not reproduce a useful analogue of Riemannian distance?
\end{problem}
\begin{solution}
Timelike proper time is maximized rather than minimized by local geodesics, while null curves have zero proper time. The causal order is an essential part of the geometry.
\end{solution}

\begin{problem}[Reverse-triangle problem]\label{prob:vii30-09}
Explain heuristically why Lorentzian time separation has a reverse triangle inequality.
\end{problem}
\begin{solution}
A maximizing proper-time curve from \(p\) to \(r\) competes with the concatenation of maximizing segments \(p\to q\) and \(q\to r\); the supremum from \(p\) to \(r\) is at least their sum.
\end{solution}

\begin{problem}[Curvature problem]\label{prob:vii30-10}
Which curvature tensors from VII/26--VII/28 remain available Lorentzianly?
\end{problem}
\begin{solution}
All of them: Riemann, Ricci, scalar, and Weyl tensors are defined using the same Levi--Civita formalism.
\end{solution}

\begin{problem}[Hyperboloid tangent problem]\label{prob:vii30-11}
Why is the tangent metric on the hyperboloid positive definite although the ambient metric is Lorentzian?
\end{problem}
\begin{solution}
The normal to the upper hyperboloid is timelike. Its Lorentz-orthogonal complement is therefore spacelike and positive definite.
\end{solution}

\begin{problem}[Comparison problem]\label{prob:vii30-12}
Summarize the principal structural difference between Riemannian and Lorentzian geometry.
\end{problem}
\begin{solution}
The differential formulas are largely shared, but positive definiteness is replaced by causal cone structure, changing length extremals, orthogonality, global separation, and admissible curves.
\end{solution}

\section{Challenge problems}
\begin{problem}[Twin-paradox geometry]\label{prob:vii30-ch01}
Use proper-time geometry to explain why a straight timelike Minkowski geodesic has greater proper time than a nearby broken timelike path with the same endpoints.
\end{problem}
\begin{solution}
For inertial segments the reverse Cauchy--Schwarz inequality implies proper-time concavity. Summing the broken segments gives less total proper time than the direct geodesic, illustrating local/global maximization in flat spacetime.
\end{solution}

\begin{problem}[Null reparametrization]\label{prob:vii30-ch02}
Why can a null geodesic have zero proper time while still possessing a meaningful affine parameter?
\end{problem}
\begin{solution}
Proper time uses \(\sqrt{-g(\dot\gamma,\dot\gamma)}\), which vanishes for null tangents. Affine parameter instead comes from the connection equation \(\nabla_{\dot\gamma}\dot\gamma=0\) and remains nontrivial.
\end{solution}

\begin{problem}[Conformal causal invariance]\label{prob:vii30-ch03}
What happens to causal type under a positive conformal rescaling \(\widetilde g=e^{2\varphi}g\)?
\end{problem}
\begin{solution}
The squared norm is multiplied by the positive factor \(e^{2\varphi}\), so its sign is unchanged. Timelike, null, and spacelike cones are therefore conformally invariant.
\end{solution}

\begin{problem}[Riemannian versus Lorentzian spheres]\label{prob:vii30-ch04}
Why does the ordinary sphere inherit a Riemannian metric from Euclidean space while the hyperboloid inherits one from Lorentzian space?
\end{problem}
\begin{solution}
In both cases the tangent space is the orthogonal complement of the normal. The Euclidean normal has positive norm and leaves a positive tangent form; the hyperboloid normal is timelike, whose Lorentz-orthogonal complement is likewise positive definite.
\end{solution}

\begin{problem}[Boundary of causal future]\label{prob:vii30-ch05}
In flat Minkowski space, why is the boundary of \(I^+(p)\) generated by null geodesics?
\end{problem}
\begin{solution}
Timelike points satisfy a strict cone inequality; the boundary occurs at equality, where the separation is null. Straight null rays are affine geodesics and generate this cone boundary.
\end{solution}

\section*{Bridge to VII/31}
VII/31 returns to positive-definite geometry and uses the Lorentzian hyperboloid as one of several equivalent models of the hyperbolic plane.
